\section{Scope, conventions, and the role of complete local subtraction}
\label{scope:section}

\subsection{Relation to the current manuscript and incoming reports}

The source baseline is the repository revision
\src{9035bb0eaaf70f3398711b489733093dcc52e946}, whose manuscript is
organized into four volumes and whose current incoming directory contains
seventeen ZIP reports \cite{repo:main}. The work here follows the
Hurwitz--Stieltjes and harmonic-identity spine. The source review included
the manuscript's structure and status records, the relevant finite-part
chapters, and the incoming reports bearing on the questions below.
It is a targeted mathematical audit, not a claim that every theorem in
the four volumes and all incoming archives has been re-proved.

There are already substantial results at this baseline. The cubic and
quartic digamma bridges, single-resolvent complex powers at nonpositive
resonances, all-degree formal Tornheim relation spaces, and the
Stieltjes-index-zero collision theorem are prior work
\cite{repo:cubic,repo:barnes,repo:quartic,repo:reflection}.
Several questions in older incoming reports have consequently become
stale. We preserve that distinction throughout: a theorem that was
already proved in a newer report is not presented as a new result here.

The concrete contributions and their limits are as follows.

\begin{longtable}{>{\raggedright\arraybackslash}p{0.24\textwidth}
                  >{\raggedright\arraybackslash}p{0.68\textwidth}}
\toprule
Target & Result in this article\\
\midrule
\endhead
Several resolvents with independent exponents &
Theorems~\ref{mr:fourier}, \ref{mr:hurwitz}, and \ref{mr:completion}
give the complete Fourier spectrum, the Hurwitz transform, and the
local completion for any finite number of nonreal parameters. Opposite
orientations are included. This answers the transformation and
normalization portions of Question 10 in \cite{repo:quartic}.\\[5pt]
Mixed logarithmic products &
Theorem~\ref{ml:allStieltjes} gives a finite expression of depth at
most two for two opposite logarithmic resolvents, at every Stieltjes
index and every ordinary argument-derivative order. It supplies a
specific instance of the mixed-product question in \cite{repo:quartic}.\\[5pt]
The fifth digamma moment &
Theorem~\ref{hp:fifth} resolves Question 1 of
\cite{repo:quartic}, with all finite corrections and a convergent
polygamma representation. It also supplies the first new case requested
by Question F3 of \cite{repo:barnes}.\\[5pt]
Arbitrary digamma powers &
Theorem~\ref{hp:allpowers} gives a terminating reduction to ordinary
convergent multiple zeta values and specified first Laurent coefficients.
Corollary~\ref{hp:primitive} fixes an antiderivative for every power.
No claim of an arithmetically minimal output basis is made.\\[5pt]
Higher-index geometric collisions &
Theorem~\ref{gcx:main} proves the total-collision limit for arbitrary
Stieltjes indices and ordinary derivatives, with an analytic local
weight and no restriction on gap ratios. This resolves the higher-index
extension posed in \cite{repo:reflection} at the level of continuity
after complete subtraction.\\[5pt]
Gaussian period conjectures &
The frozen $S_6$ and revised $S_8$ remain open. Neither a new period
certificate nor a counterexample is obtained. The earlier rejected
$S_8$ vector and the stated formal relation obstructions retain their
separate established status \cite{repo:main,repo:exact}.\\
\bottomrule
\end{longtable}

The word \emph{new} refers here to progress relative to this
pinned source corpus. Classical beta integrals, Gamma expansions,
harmonic products, and multiple-zeta dualities are cited as such.
The external literature check was focused; it does not establish global
priority for every consequence or specialization.

\subsection{Special-function and coefficient conventions}

The generalized Stieltjes functions are normalized by
\begin{equation}\label{scope:stieltjes}
 \zeta(1+u,x)=\frac1u+
       \sum_{m\ge0}\frac{(-1)^m}{m!}\gamma_m(x)u^m,
 \qquad \gamma_m=\gamma_m(1),\quad \gamma=\gamma_0.
\end{equation}
Thus $\gamma_0(x)=-\psi(x)$, where $\psi=\Gamma'/\Gamma$.
Equivalently,
\[
 \zeta(1-s,x)+\frac1s=\sum_{m\ge0}\gamma_m(x)\frac{s^m}{m!}.
\]
A superscript $(p)$ on $\gamma_m$ or $\psi$ means an ordinary derivative
in its argument. A prime on a zeta function of a single variable denotes
an order derivative unless an argument is explicitly specified.
The standard Hurwitz recurrence and differentiation formulas are
\cite{dlmf:hurwitz,dlmf:polygamma}
\begin{align}
 \gamma_m(x)&=\frac{\log^m x}{x}+\gamma_m(1+x),
       \label{scope:recurrence}\\
 \partial_x^p\zeta(s,x)&=(-1)^p(s)_p\zeta(s+p,x).
       \label{scope:hurwitzderivative}
\end{align}
Here $(s)_p=s(s+1)\cdots(s+p-1)$ is the rising factorial.
The symbol $[u^j]F$ means a Taylor or Laurent coefficient, including
no factorial. A derivative is recovered by multiplying a Taylor
coefficient by the appropriate factorial. This distinction matters at
every harmonic Laurent coefficient used below.

Our multiple zeta convention is strict and decreasing:
\[
 Z_d(s_1,\ldots,s_d)
   =\sum_{n_1>\cdots>n_d\ge1}
       n_1^{-s_1}\cdots n_d^{-s_d},
\]
in its convergence region, with meromorphic continuation where stated.
The notation $\zeta(s,t)$ in the polylogarithmic section denotes the
same depth-two function. For a nonnegative integer $N$,
$H_N^{(r)}=\sum_{n=1}^N n^{-r}$; multiple harmonic words use the same
strict ordering. No weakly ordered sum is silently substituted for a
strict one.

\subsection{Coordinate finite parts}

All endpoint finite parts are taken in the original integration
coordinate. For an integrand with a finite power--log singular expansion,
integrate after removing a small endpoint interval, expand the result
in the cutoff, and retain its constant term. At two distinct endpoints
the cutoffs are independent. In particular,
\begin{align}
 \FP\int_0^1 x^{\alpha-1}\log^j x\,\dd x
   &=\frac{(-1)^j j!}{\alpha^{j+1}}
                    &&(\alpha\ne0),\label{scope:monomial}\\
 \FP\int_0^1 x^{-1}\log^j x\,\dd x&=0
                    &&(j\ge0).\label{scope:resonantmonomial}
\end{align}
The first equality follows by meromorphic continuation from
$\Re\alpha>0$, or directly from the cutoff primitive.
At $\alpha=0$ the second equality is the coordinate constant;
it is not a value obtained by inserting zero into the first formula.

For periodic Stieltjes products, the local singularity occurs on the
right of each seam. We remove a right interval of ambient length
$\eta$ at every singular point and take independent constant terms.
Ordinary derivatives on the open intervals are used throughout.
Differentiating a periodic distribution that has already been extended
across a seam is a different operation and can carry additional
distributional contacts. Section~\ref{gcx:section} specifies its
convention explicitly.

For fixed distinct singular points, independent finite parts are linear
and commute with fixed parameter derivatives once enough local
monomials have been subtracted uniformly. This elementary fact is the
common analytic step behind the spectral and geometric constructions.
It does not authorize taking an arbitrary path through intersecting
parameter poles.

\begin{lemma}[Full-amplitude completion]\label{scope:completionlemma}
Let a parameter-dependent integrand near an endpoint have a finite
expansion whose singular terms are
$A_j(\mathbf t)x^{\alpha_j(\mathbf t)-1}$, with holomorphic amplitudes
and affine exponents, none identically zero. Suppose the subtracted remainder and each fixed
parameter derivative are uniformly integrable near $\mathbf t=0$.
For every $j$ such that $\alpha_j(0)=0$, subtract
$A_j(\mathbf t)/\alpha_j(\mathbf t)$ from its meromorphically continued
integral. The resulting local germ is holomorphic, and its Taylor
coefficients are exactly the coordinate finite parts of the
corresponding derivatives of the integrand. Power--log monomials follow
by differentiating in an auxiliary exponent.
\end{lemma}

\begin{proof}
On a fixed endpoint interval $(0,\rho)$, the integrated monomial is
\[
 A_j(\mathbf t)\frac{\rho^{\alpha_j(\mathbf t)}}
                            {\alpha_j(\mathbf t)}
 =\frac{A_j(\mathbf t)}{\alpha_j(\mathbf t)}
  +A_j(\mathbf t)
       \frac{\rho^{\alpha_j(\mathbf t)}-1}{\alpha_j(\mathbf t)}.
\]
The second term is entire in $\alpha_j$ and generates the cutoff
constants of all its logarithmic derivatives. Nonresonant monomials are
already holomorphic at the base point. Uniform integrability justifies
coefficient extraction from the remainder. The artificial splitting
point cancels against the rest of the integral.
\end{proof}

The amplitude in this lemma is the complete function
$A_j(\mathbf t)$, not merely its restriction to the polar hyperplane.
The latter specifies a residue but omits regular coefficients that
contribute to the prescribed finite part. The two-orientation example
in Section~\ref{mr:section} makes this distinction exact: different
uncompleted parameter paths give different finite limits for the same
specialized function.
